\chapter{计算预算感知的训练数据选择}

\section{章节引言}

第二章指出，现���训练数据选择方法隐含假设最优数据子集独立于计算预算。本章从一个直觉出发挑战这一假设：在只能训练10轮和可以训练100轮时，"最有价值的数据"应当是相同的吗？

直觉上，答案是否定的。当计算预算有限时，模型仅有时间学习数据中的简单模式\cite{rahaman2018spectral}，此时复杂或含噪声的样本不仅无法被有效利用，反而可能干扰学习过程；当计算预算充裕时，模型有能力从多样化甚至含噪声的数据中提取更丰富的信息，此时仅使用简单数据反而限制了模型的最终能力。这意味着最优数据子集应当依赖于可用的计算预算。

为验证这一直觉，本文设计了一组探索性实验。在由不同噪声水平构成的异构多源数据上，固定数据选择策略（如始终选择最干净的数据源，或始终均匀混合所有源），然后在不同计算预算下评估各策略的表现。图~\ref{fig:cads-exploratory}展示了实验结果。% [VERIFY: 实验设置——CIFAR-10, 5个噪声水平的数据源]

% [图3-1: 不同计算预算下低频与高频数据训练的验证损失对比]
\begin{figure}[htb]
  \centering
  % \includegraphics[width=0.8\textwidth]{figures/cads-exploratory.pdf}
  \fbox{\parbox{0.8\textwidth}{\centering\vspace{2.5cm}
    \textbf{占位图：不同计算预算下低频与高频数据训练的验证损失对比}\\[0.5em]
    对应原论文Fig.1，需手绘
  \vspace{2.5cm}}}
  \caption{不同计算预算下低频与高频数据训练的验证损失对比}
  \label{fig:cads-exploratory}
\end{figure}

实验结果清晰地展示了两个关键现象：第一，没有任何固定数据选择策略能在所有计算预算下一致领先——低预算下偏好干净数据的策略表现最优，而高预算下纳入更多数据源的策略反而更优；% [VERIFY: 具体数字]
第二，不同计算预算下的最优数据配比截然不同。这些发现指向一个此前被忽视的结论：\textbf{最优数据选择策略是计算预算的函数}，预算不应被隐含假设为充裕，而应作为数据选择的显式输入。

基于上述观察，本章提出CADS（Computation-Aware Data Selection）框架，其核心思想是将计算预算$C$作为双层优化中的显式约束，让数据选择策略自动适应给定预算。图~\ref{fig:cads-framework}展示了CADS框架的总体架构。

% [图3-2: CADS框架总览图]
\begin{figure}[htb]
  \centering
  % \includegraphics[width=0.9\textwidth]{figures/cads-framework.pdf}
  \fbox{\parbox{0.9\textwidth}{\centering\vspace{2.5cm}
    \textbf{占位图：CADS框架总览图}\\[0.5em]
    展示双层优化结构、惩罚松弛、CADS-E/S两种变体的关系\\
    需手绘
  \vspace{2.5cm}}}
  \caption{CADS框架总览——双层优化结构、惩罚松弛与两种变体}
  \label{fig:cads-framework}
\end{figure}

实现这一思想面临两个关键技术挑战：（1）预算约束导致内层训练无法收敛，使得标准双层优化方法（如基于隐式微分的方法）失效；（2）数据选择过程本身不应消耗过多计算预算，否则将违背"高效利用数据"的初衷。本章将依次介绍CADS如何应对这两个挑战。


\section{问题定义}

\subsection{计算约束下的数据选择建模}

考虑训练集$\mathcal{D}_{\text{trn}} = \{(x_i, y_i)\}_{i=1}^{N}$和验证集$\mathcal{D}_{\text{val}}$。计算预算$C$以前向-反向传播的总次数度量。数据选择掩码$\mathbf{m} \in \{0,1\}^N$指示哪些样本被选入训练子集。

在计算预算$C$的约束下，数据选择问题可建模为如下双层优化：
\begin{equation}\label{eq:cads-bilevel}
  \min_{\mathbf{s}} \mathbb{E}_{\mathbf{m} \sim p(\mathbf{m}|\mathbf{s})}[\Lval(\btheta_C(\mathbf{m}))]
\end{equation}
其中$\btheta_C(\mathbf{m})$表示在子集$\mathbf{m}$上训练恰好$C$步后得到的模型参数，$p(\mathbf{m}|\mathbf{s})$是以$\mathbf{s}$为参数的采样分布。外层目标是找到使验证损失最小的选择策略$\mathbf{s}$。

这一建模与标准核心子集选择问题\cite{zhou2022probabilistic}的关键区别在于对$\btheta_C(\mathbf{m})$的定义：标准方法假设内层训练到最优$\btheta^*(\mathbf{m})$，而CADS显式将内层训练限制在$C$步。当$C \to \infty$时，CADS退化为标准核心子集选择问题。

\subsection{预算约束带来的根本困难}

标准双层优化假设内层求解到最优$\btheta^*(\mathbf{m})$，在此条件下可利用隐式函数定理\cite{domke2012generic, grazzi2020iteration}计算外层梯度。然而，预算约束下内层仅运行$C$步，$\btheta_C(\mathbf{m}) \neq \btheta^*(\mathbf{m})$，内层最优性条件不成立，隐式微分的理论前提被破坏。

一种替代方案是迭代微分\cite{maclaurin2015gradient, shaban2019truncated}——对$C$步训练过程做截断反向传播。然而，这要求存储整个训练轨迹的中间状态，内存和计算开销均与$C$成正比\cite{chen2016training}，对大预算场景不可行。

因此，CADS需要一种不依赖内层最优性、也不需要对训练过程做梯度回传的新方法来估计外层梯度。


\section{方法：CADS框架}

\subsection{双层优化建模与策略梯度基线}

\textbf{概率化重参数化。}CADS首先将离散掩码$\mathbf{m}$的组合优化问题转化为连续参数$\mathbf{s}$的优化问题。定义参数化采样分布$p(\mathbf{m}|\mathbf{s})$，外层目标变为关于$\mathbf{s}$的期望优化。这一重参数化将搜索空间从$2^N$个候选子集降低为对连续参数$\mathbf{s}$的梯度优化。

\textbf{策略梯度估计。}利用REINFORCE策略梯度\cite{zhou2022probabilistic}，外层梯度可以表示为：
\begin{equation}\label{eq:cads-reinforce}
  \nabla_{\mathbf{s}}\mathbb{E}[\Lval] = \mathbb{E}_{\mathbf{m} \sim p(\mathbf{m}|\mathbf{s})}[\Lval(\btheta_C(\mathbf{m})) \cdot \nabla_{\mathbf{s}} \log p(\mathbf{m}|\mathbf{s})]
\end{equation}
这一估计具有关键优势：它无需对内层训练过程做梯度回传（即无需Hessian计算），外层梯度仅需评估验证损失$\Lval$和采样分布的对数梯度$\nabla_{\mathbf{s}} \log p(\mathbf{m}|\mathbf{s})$。从强化学习的视角看，数据选择$\mathbf{m}$是"动作"，验证损失是"奖励"，$\mathbf{s}$是"策略参数"。

\textbf{Bilevel-CADS的局限。}基于策略梯度的方法（称为Bilevel-CADS）在概念上验证了预算感知数据选择的有效性，但其计算开销仍然很高：每次外层更新需要$K$次独立的完整训练——采样$K$个不同的$\mathbf{m}$，并各自训练$C$步——总计算量为$K \times C$。即使$K$仅取5，在中等预算下也需要5倍完整训练的开销，实际应用中难以接受。Bilevel-CADS作为概念验证，证明了预算感知数据选择的价值，但距离实用化尚有差距。

\subsection{惩罚松弛与单层目标}

为克服Bilevel-CADS的效率瓶颈，CADS提出将双层优化转化为单层优化的惩罚松弛方法\cite{ye2022bome, kwon2023fully}。

\textbf{核心思想。}引入惩罚项，构造如下单层目标：
\begin{equation}\label{eq:cads-penalty}
  \mathcal{L}^{\alpha}_{\text{CADS}} = \mathbb{E}_{p(\mathbf{m}|\mathbf{s})}[\Lval(\btheta) + \alpha(\Ltrn(\btheta, \mathbf{m}) - l(|\mathbf{m}|))^2]
\end{equation}
其中$\alpha$为惩罚系数，$l(|\mathbf{m}|)$为"在大小为$|\mathbf{m}|$的子集上训练$C$步后的最小可达训练损失"。惩罚项的直觉含义是：模型在当前子集上的训练损失应接近该子集大小所对应的"目标损失"——即模型应充分利用所选数据，既不欠拟合（训练损失远高于$l(|\mathbf{m}|)$）也不过拟合（训练损失远低于$l(|\mathbf{m}|)$）。

这一松弛将双层嵌套结构消解为单层联合优化问题：模型参数$\btheta$和选择策略参数$\mathbf{s}$可通过标准SGD交替更新，无需内外层迭代。

\textbf{一维损失近似$l(|\mathbf{m}|)$。}惩罚项中的$l(|\mathbf{m}|)$——即在给定子集大小下$C$步训练后的最小可达损失——如果直接求解，需要对每个候选子集大小做完整训练，计算上不可行。

CADS的一个关键发现是：$l(|\mathbf{m}|)$关于子集大小$|\mathbf{m}|$近似呈对数线性关系（log-linear）。这意味着只需在3至5个不同子集大小上做短训练、记录最终训练损失，即可拟合一条对数线性曲线来近似$l(\cdot)$。% [VERIFY: 拟合细节和精度]
图~\ref{fig:cads-loss-approx}展示了这一近似的效果。

% [图3-3: 一维损失近似的拟合效果]
\begin{figure}[htb]
  \centering
  % \includegraphics[width=0.6\textwidth]{figures/cads-loss-approx.pdf}
  \fbox{\parbox{0.6\textwidth}{\centering\vspace{2cm}
    \textbf{占位图：一维损失近似的拟合效果}\\[0.5em]
    不同子集大小下的实际训练损失与拟合曲线对比\\
    对应原论文Fig.6/附录，需手绘
  \vspace{2cm}}}
  \caption{一维损失近似——不同子集大小下的实际训练损失与对数线性拟合曲线}
  \label{fig:cads-loss-approx}
\end{figure}

这一近似将惩罚项的计算从"每个候选子集独立训练"降低为"一次性拟合 + 查表"，是CADS效率提升的关键所在。

\textbf{效率分析。}CADS的总额外计算开销为$(5/A + 2\gamma)C$，其中$A$为摊销因子（一次选择可复用于多次训练���，$\gamma < 1$为验证集评估比例。% [VERIFY: 具体公式]
作为对比，表~\ref{tab:cads-cost}将CADS与主流数据选择方法的计算开销进行了比较：基于训练动态的方法（如Forgetting\cite{toneva2018empirical}、GraNd/EL2N\cite{paul2021deep}）需要$1.2C$至$2C$的额外开销，而更精细的方法（如Influence Functions、Data Shapley）需要$10C$至$101C$以上的开销。% [VERIFY: Table 6数据]
CADS在大预算场景下可低至约$2\gamma C$的额外开销，且其摊销特性意味着一次选择可复用于多次独立训练。

\subsection{CADS-E：样本级选择}

CADS-E是CADS框架的样本级变体，为每个训练样本独立建模选择概率。

在采样分布设计上，每个样本$i$对应一个独立的Bernoulli参数$s_i \in [0,1]$，表示该样本被选入子集的概率。采样分布为���
\begin{equation}\label{eq:cads-e}
  p(\mathbf{m}|\mathbf{s}) = \prod_{i=1}^{N} s_i^{m_i}(1-s_i)^{1-m_i}
\end{equation}
对数梯度$\partial \log p / \partial s_i = m_i/s_i - (1-m_i)/(1-s_i)$，计算简单且可并行。

CADS-E的优势在于细粒度控制：每个样本拥有独立的选择概率，可识别出特定的噪声样本或特别有价值的样本。其局限在于参数量等于训练集大小$N$，对大规模数据集存在优化效率问题。因此，CADS-E更适用于中小规模数据集或需要样本级精细控制的场景。

\subsection{CADS-S：来源级选择}

CADS-S是CADS框架的来源级变体，将训练数据按来源分组，在来源级别而非样本级别进行选择。

假设训练数据由$n$个来源$\{\mathcal{D}_j\}_{j=1}^{n}$组成，每个来源$j$的采样比例$r_j$服从截断高斯分布：
\begin{equation}\label{eq:cads-s}
  r_j \sim \mathcal{N}(s_j, \sigma_t^2), \quad r_j \in [0,1]
\end{equation}
其中$s_j$为可学习的中心参数，$\sigma_t = 0.99^t \sigma_0$为随训练步数$t$递减的方差。方差退火机制使得选择策略在训练初期保持充分探索，在后期逐渐稳定。

CADS-S的参数量仅为来源数$n$（通常$n \ll N$），适用于多源异构数据的场景，如不同质量的数据集混合、不同域的数据配比。在这类场景中，关键决策是各来源的混合比例，而非逐个样本的取舍。

CADS-S揭示了CADS框架最重要的发现之一：数据选择策略随计算预算系统性地变化。在CIFAR-10异构数据源实验中，低预算时CADS-S几乎只选择最干净的数据源，高噪声源的采样权重趋近于零；随着预算增加，CADS-S逐步纳入中等噪声水平的数据源；极高噪声源（噪声率$>$60\%）在所有预算下均被排除。% [VERIFY: 具体阈值]
这一行为直接体现了数据选择与计算预算的交互：预算决定了模型能"消化"多大难度的数据。


\section{实验}

\subsection{实验设置}

本章设计了四组实验，覆盖不同规模和模态的场景：MNIST\cite{lecun1998mnist}用于小规模概念验证，CIFAR-10\cite{krizhevsky2009cifar}用于异构数据源的系统评估（构造5个不同噪声水平的数据源），DomainNet\cite{domainnet}用于大规模多域场景（6个视觉域），GPT-2\cite{radford2019language}指令微调用于语言模态验证。% [VERIFY: 各实验的详细设置]

基线方法包括：随机选择、均匀混合（Uniform）、最优单源（Best-source）、概率双层核心子集选择（PBCS\cite{zhou2022probabilistic}）以及全数据训练。评价指标为分类准确率（视觉任务）和困惑度（语言任务）。

\subsection{小规模验证}

首先在MNIST数据集上验证Bilevel-CADS基线的概念有效性。表~\ref{tab:cads-mnist}展示了不同数据选择方法在不同初始子集大小下的准确率对比。

% [表3-1: MNIST上不同数据选择方法在不同初始子集大小下的准确率对比]
\begin{table}[htb]
  \centering
  \caption{MNIST上不同数据选择方法在不同初始子集大小下的准确率对比}
  \label{tab:cads-mnist}
  \fbox{\parbox{0.8\textwidth}{\centering\vspace{1.5cm}
    \textbf{占位表：对应原论文Table 1}\\[0.5em]
    % [VERIFY: Table 1]
  \vspace{1.5cm}}}
\end{table}

Bilevel-CADS在不同初始子集大小（200、400、600、800）下均优于随机选择和PBCS，平均准确率达到92.80\%。% [VERIFY: Table 1]
图~\ref{fig:cads-convergence}展示了Bilevel-CADS优化过程中子集大小的收敛轨迹——无论从何种初始大小出发，子集大小均收敛到相近的最终值，表明优化过程对初始化不敏感。

% [图3-4: Bilevel-CADS在不同初始子集大小下的子集大小收敛过程]
\begin{figure}[htb]
  \centering
  % \includegraphics[width=0.6\textwidth]{figures/cads-convergence.pdf}
  \fbox{\parbox{0.6\textwidth}{\centering\vspace{2cm}
    \textbf{占位图：子集大小收敛过程}\\[0.5em]
    对应原论文Fig.2，需手绘
  \vspace{2cm}}}
  \caption{Bilevel-CADS在不同初始子集大小下的子集大小收敛过程}
  \label{fig:cads-convergence}
\end{figure}

这一实验的意义在于证明预算感知的数据选择确实能提升模型性能，为后续高效方法的开发提供了动机基础。

\subsection{异构数据源实验}

\textbf{CIFAR-10实验。}在CIFAR-10上构造由5个不同噪声水平（0\%至90\%）的数据源组成的异构数据集，计算预算从90,000到225,000样本不等。表~\ref{tab:cads-cifar}展示了各方法在不同预算下的准确率。

% [表3-2: CIFAR-10异构数据源实验]
\begin{table}[htb]
  \centering
  \caption{CIFAR-10异构数据源实验——不同计算预算下各方法的准确率对比}
  \label{tab:cads-cifar}
  \fbox{\parbox{0.8\textwidth}{\centering\vspace{1.5cm}
    \textbf{占位表：对应原论文Table 2}\\[0.5em]
    % [VERIFY: Table 2]
  \vspace{1.5cm}}}
\end{table}

CADS-S平均准确率63.08\%，显著优于Best-source基线的61.65\%和全数据训练。% [VERIFY: Table 2]
更重要的是，CADS-S在所有计算预算下一致领先，而其他方法的相对优劣随预算变化。图~\ref{fig:cads-cifar-ratio}展示了CADS-S在四种不同预算下的采样比例演化过程，清晰地展示了选择策略如何随预算自适应调整。

% [图3-5: CIFAR-10上不同计算预算下CADS-S的采样比例演化过程]
\begin{figure}[htb]
  \centering
  % \includegraphics[width=\textwidth]{figures/cads-cifar-ratio.pdf}
  \fbox{\parbox{\textwidth}{\centering\vspace{2cm}
    \textbf{占位图：CIFAR-10采样比例演化过程}\\[0.5em]
    四个子图分别对应四种预算，对应原论文Fig.3，需手绘
  \vspace{2cm}}}
  \caption{CIFAR-10上不同计算预算下CADS-S的采样比例演化过程}
  \label{fig:cads-cifar-ratio}
\end{figure}

\textbf{DomainNet实验。}在包含6个视觉域的大规模DomainNet数据集上，CADS-S进一步验证了其可扩展性。表~\ref{tab:cads-domainnet}展示了结果。

% [表3-3: DomainNet上不同初始采样比例下各方法的准确率对比]
\begin{table}[htb]
  \centering
  \caption{DomainNet上不同初始采样比例下各方法的准确率对比}
  \label{tab:cads-domainnet}
  \fbox{\parbox{0.8\textwidth}{\centering\vspace{1.5cm}
    \textbf{占位表：对应原论文Table 3}\\[0.5em]
    % [VERIFY: Table 3]
  \vspace{1.5cm}}}
\end{table}

CADS-S平均准确率42.27\%，远优于均匀采样基线的36.22\%，最高提升达14.42\%。% [VERIFY: Table 3]
图~\ref{fig:cads-domainnet-ratio}展示了DomainNet上的采样比例演化。

% [图3-6: DomainNet上CADS-S的采样比例演化过程]
\begin{figure}[htb]
  \centering
  % \includegraphics[width=0.8\textwidth]{figures/cads-domainnet-ratio.pdf}
  \fbox{\parbox{0.8\textwidth}{\centering\vspace{2cm}
    \textbf{占位图：DomainNet采样比例演化过程}\\[0.5em]
    对应原论文Fig.4，需手绘
  \vspace{2cm}}}
  \caption{DomainNet上CADS-S的采样比例演化过程}
  \label{fig:cads-domainnet-ratio}
\end{figure}

这两组实验共同验证了核心发现：CADS-S能自动学习到适配当前计算预算的数据配比策略，在所有预算设定下一致优于固定策略基线。

\subsection{指令微调实验}

为验证CADS在语言模态上的通用性，本文在GPT-2\cite{radford2019language}指令微调任务上进行了实验。

\textbf{双源实验。}使用Alpaca-GPT4\cite{peng2023instruction}（1K条高质量指令）和Alpaca\cite{alpaca}（9K条标准质量指令）作为两个数据源，计算预算从$1\times10^4$到$5\times10^4$。表~\ref{tab:cads-alpaca}展示了各方法在不同预算下的困惑度对比。

% [表3-4: GPT-2指令微调——不同计算预算下各方法的困惑度对比]
\begin{table}[htb]
  \centering
  \caption{GPT-2指令微调——不同计算预算下各方法的困惑度对比}
  \label{tab:cads-alpaca}
  \fbox{\parbox{0.8\textwidth}{\centering\vspace{1.5cm}
    \textbf{占位表：对应原论文Table 4}\\[0.5em]
    % [VERIFY: Table 4]
  \vspace{1.5cm}}}
\end{table}

CADS-S在困惑度指标上一致优于基线方法。% [VERIFY: Table 4]
图~\ref{fig:cads-alpaca-ratio}展示了两个数据源的采样比例随预算的动态变化，再次体现了预算对最优数据配比的影响。

% [图3-7: Alpaca数据集上不同计算预算下的采样比例动态]
\begin{figure}[htb]
  \centering
  % \includegraphics[width=0.7\textwidth]{figures/cads-alpaca-ratio.pdf}
  \fbox{\parbox{0.7\textwidth}{\centering\vspace{2cm}
    \textbf{占位图：Alpaca采样比例动态}\\[0.5em]
    对应原论文Fig.5，需手绘
  \vspace{2cm}}}
  \caption{Alpaca数据集上不同计算预算下的采样比例动态}
  \label{fig:cads-alpaca-ratio}
\end{figure}

\textbf{大规模数据混合实验。}进一步在13个指令微调数据集\cite{alpaca, peng2023instruction, lian2023slimorca, openorcagpt4, gpteacher, khashabi2020unifiedqa, liu2019multi, wei2021finetuned, bukharin2023data, raffel2019exploring}上验证CADS-S的可扩展性。表~\ref{tab:cads-13source}展示了结果。

% [表3-5: 13个指令微调数据集的大规模数据混合实验结果]
\begin{table}[htb]
  \centering
  \caption{13个指令微调数据集的大规模数据混合实验结果}
  \label{tab:cads-13source}
  \fbox{\parbox{0.8\textwidth}{\centering\vspace{1.5cm}
    \textbf{占位表：对应原论文Table 5}\\[0.5em]
    % [VERIFY: Table 5]
  \vspace{1.5cm}}}
\end{table}

在13个数据源的混合场��中，CADS-S仍然能够有效地优化数据配比，证明了其在大规模多源数据选择中的实用价值。% [VERIFY: Table 5]

\subsection{计算开销分析}

表~\ref{tab:cads-cost}将CADS与主流数据选择方法的计算开销进行了系统对比。

% [表3-6: 各数据选择方法的计算开销对比分析]
\begin{table}[htb]
  \centering
  \caption{各数据选择方法的计算开销对比分析}
  \label{tab:cads-cost}
  \fbox{\parbox{0.8\textwidth}{\centering\vspace{1.5cm}
    \textbf{占位表：对应原论文Table 6}\\[0.5em]
    % [VERIFY: Table 6]
  \vspace{1.5cm}}}
\end{table}

对比结果表明，CADS的计算开销随预算规模增大而摊薄。% [VERIFY: Table 6]
在大预算场景下，CADS的额外开销远低于基于训练动态的方法（Forgetting、GraNd、EL2N\cite{toneva2018empirical, paul2021deep}），更远低于基于数据价值估计的方法（Influence Functions\cite{pruthi2020estimating}、Data Shapley）。这一效率优势使CADS成为大规模数据选择场景的实用方案。


\section{讨论}

\subsection{损失预测器的角色与局限性}

CADS的惩罚松弛方法依赖一维损失近似$l(|\mathbf{m}|)$的准确性。如果该近似偏差过大，惩罚项将引导模型进入错误的训练状态，导致选择策略失效。

对数线性近似在以下情况下可能不够准确：数据分布极端不均匀、计算预算范围跨越多个数量级、数据来源间质量差异极大。在这些极端场景下，可能需要更灵活的拟合方法，如分段线性拟合或使用Scaling Law中常用的四参数函数族\cite{kaplan2020scaling, hoffmann2022training}。

\subsection{与大模型数据混合领域的联系}

CADS-S的来源级选择与LLM预训练中的数据混合（Data Mixture）问题在形式上高度一致：二者都需要确定不同数据来源的最优混合比例\cite{albalak2024survey}。LLM领域的损失预测器和Scaling Law方法（如Chinchilla\cite{hoffmann2022training}）通过预测不同配比下的损失来指导数据混合，与CADS中的一维损失近似$l(|\mathbf{m}|)$扮演类似角色。

二者的关键异同在于：CADS显式引入计算预算作为变量，使得选择策略随预算自适应变化；而大多数LLM数据混合工作假设固定的预算（固定token数）。CADS的预算感知思路可能对LLM领域的数据配比研究提供新的视角。


\section{章节小结}

本章提出了CADS框架，首次将计算预算纳入数据选择的优化框架，揭示了一个被广泛忽视的交互：\textbf{最优数据子集是计算预算的函数}——低预算偏好简单干净的数据，高预算受益于多样化的数据。在方法层面，CADS通过概率化重参数化的策略梯度估计绕过了内层非收敛的困难，进一步通过惩罚松弛和一维损失近似将昂贵的双层优化转化为高效的单层求解。在实证层面，CADS跨视觉分类和语言模型微调任务验证了预算感知数据选择的一致优势，最高提升14.42\%。% [VERIFY]

CADS在优化数据选择时假设模型结构固定——模型作为"黑盒"接收所选数据并返回验证损失。然而，模型本身也存在冗余参数\cite{han2015learning, frankle2018lottery}，而模型容量（参数量）与数据量之间存在已知的Scaling关系\cite{kaplan2020scaling, hoffmann2022training}。一个自然的追问是：当模型容量也成为可优化变量时——例如通过剪枝压缩模型——数据选择与参数冗余之间是否存在类似的交互？第四章将回答这一问题。
